\chapter{A Empresa BNDES}

\section{Introdução}

O trabalho tem como objetivo apresentar algumas das principais características do Banco Nacional de Desenvolvimento Econômico e Social (BNDES), uma empresa pública federal brasileira. Serão abordados seu ramo de atividade, seu porte de acordo com o critério utilizado pelo próprio BNDES, além de sua missão, visão e valores.

O BNDES foi fundado em 20 de junho de 1952 e é vinculado ao Ministério do Desenvolvimento, Indústria, Comércio e Serviços. O banco tem como principal função apoiar o desenvolvimento econômico, social e ambiental do Brasil, por meio de financiamentos, investimentos e outras formas de apoio.

\section{Ramo de Atividade}

O BNDES atua no setor financeiro, especificamente como um banco de desenvolvimento. Sua atuação é diferente da de um banco comercial tradicional, pois seu principal objetivo é apoiar projetos que contribuam para o desenvolvimento do país. O banco oferece apoio financeiro para empresas, empreendedores e projetos de diferentes setores da economia. Entre suas formas de atuação estão os financiamentos, a concessão de recursos não reembolsáveis, a subscrição de valores mobiliários e a participação em fundos de investimento. O BNDES também busca incentivar áreas como inovação, desenvolvimento regional, sustentabilidade, geração de empregos e inclusão social \cite{bndes2026quemsomos}.

% \noindent\textbf{Fonte:} BNDES -- Quem Somos.

\section{Porte da Empresa Segundo o Critério BNDES}

O BNDES classifica o porte das empresas de acordo com a Receita Operacional Bruta (ROB) anual da empresa ou do grupo econômico ao qual ela pertence. De acordo com o critério utilizado atualmente pelo BNDES, a classificação é apresentada no Quadro~\ref{tab:porte} \cite{bndes2026faq}.

\begin{table}[htbp]
    \centering
    \caption{Classificação de porte de clientes segundo critério BNDES.}
    \label{tab:porte}
    \vspace{0.3cm}
    \begin{tabularx}{\textwidth}{l X}
        \toprule
        \textbf{Classificação} & \textbf{Receita Operacional Bruta Anual}                     \\
        \midrule
        Microempresa           & Menor ou igual a R\$ 360 mil                                 \\
        Pequena empresa        & Maior que R\$ 360 mil e menor ou igual a R\$ 4,8 milhões     \\
        Média empresa          & Maior que R\$ 4,8 milhões e menor ou igual a R\$ 300 milhões \\
        Grande empresa         & Maior que R\$ 300 milhões                                    \\
        \bottomrule
    \end{tabularx}
    \vspace{0.2cm}
    \begin{flushleft}
        \centering
        \footnotesize\textbf{Fonte:} BNDES -- Perguntas Frequentes: classificação de porte das empresas.
    \end{flushleft}
\end{table}

É importante destacar que essa classificação é utilizada pelo BNDES principalmente para definir o porte de seus clientes e direcionar suas condições de apoio financeiro. O BNDES, especificamente, é uma empresa pública federal, e não deve ser classificado simplesmente como uma \enquote{grande empresa} utilizando essa tabela de clientes.

% \noindent\textbf{Fonte:} BNDES -- Perguntas Frequentes: classificação de porte das empresas.

\section{Missão}

Conforme a página oficial da instituição, a missão atual do BNDES é: \enquote{Retomar o protagonismo do BNDES no desenvolvimento econômico, social e ambiental brasileiro.}\cite{bndes2026planejamento}

A missão demonstra que a instituição busca exercer um papel importante no desenvolvimento do Brasil, considerando não apenas o crescimento econômico, mas também questões sociais e ambientais \cite{bndes2026planejamento}.

% \noindent\textbf{Fonte:} BNDES -- Planejamento Estratégico.

\section{Visão}

A visão atual do BNDES consiste em \enquote{Ser um banco de desenvolvimento sustentável, digital, inclusivo, inovador, industrializante e tecnológico.} \cite{bndes2026planejamento}

A visão demonstra onde a instituição pretende chegar no futuro. Ela destaca a importância da sustentabilidade, da tecnologia, da inovação e da inclusão em sua atuação \cite{bndes2026planejamento}.

% \noindent\textbf{Fonte:} BNDES -- Planejamento Estratégico.

\section{Valores}

Os valores institucionais atuais do BNDES são \cite{bndes2026planejamento}:

\begin{itemize}
    \item \textbf{Compromisso com desenvolvimento:} representa o compromisso da instituição em contribuir para o desenvolvimento econômico, social e ambiental do país.
    \item \textbf{Espírito público:} está relacionado à atuação do BNDES em benefício do interesse público e do desenvolvimento da sociedade brasileira.
    \item \textbf{Ética:} refere-se à importância de agir com responsabilidade, integridade, transparência e respeito nas atividades realizadas pela instituição.
    \item \textbf{Excelência:} está relacionada à busca pela qualidade e pelo aprimoramento contínuo na atuação do banco.
\end{itemize}

% \noindent\textbf{Fonte:} BNDES -- Planejamento Estratégico.

\section{Considerações}

A análise do BNDES permite compreender a importância de uma instituição financeira voltada ao desenvolvimento do país. Como empresa pública federal, o banco possui um papel relevante no apoio a investimentos e projetos que podem contribuir para o crescimento econômico e para o desenvolvimento social e ambiental do Brasil.

Seu ramo de atividade está relacionado ao setor financeiro e ao financiamento de projetos de desenvolvimento. O BNDES utiliza a Receita Operacional Bruta para classificar o porte de seus clientes, enquanto sua própria natureza jurídica é de empresa pública federal.

Sua missão, visão e valores demonstram uma preocupação com desenvolvimento sustentável, inovação, inclusão, ética e excelência. Dessa forma, o BNDES procura alinhar sua atuação financeira com objetivos econômicos, sociais e ambientais.